% Part: applied-modal-logics
% Chapter: epistemic-logic
% Section: language-epistemic-logic

\documentclass[../../../include/open-logic-section]{subfiles}

\begin{document}

\olfileid[pt]{aml}{el}{pal}

\olsection{Lógica de anúncios públicos}

As lógicas epistêmicas dinâmicas permitem representar como o
conhecimento dos agentes muda ao longo do tempo ou à medida que eles
obtêm novas informações. Muitas delas representam mudanças no
conhecimento por meio de \emph{eventos} ou \emph{atualizações}
informacionais. O tipo mais básico de atualização é um anúncio público,
no qual alguma fórmula é anunciada de modo verídico e todos os agentes
presenciam juntos o acontecimento. Para isso, ampliamos a linguagem da
seguinte maneira.

\begin{defn}
Seja $G$ um conjunto de símbolos de agentes. A linguagem básica da
lógica epistêmica multiagente com anúncios públicos contém:
\begin{enumerate}
  \tagitem{prvFalse}{A constante proposicional de !!{falsity}~$\lfalse$.}{}
  \tagitem{prvTrue}{A constante proposicional de !!{truth}~$\ltrue$.}{}
  \item Um conjunto de elementos !!{denumerable}s, as !!{propositional variable}s: $\Obj
    p_0$, $\Obj p_1$, $\Obj p_2$, \dots
  \item Os conectivos proposicionais: \startycommalist
  \iftag{prvNot}{\ycomma $\lnot$ (negação)}{}%
  \iftag{prvAnd}{\ycomma $\land$ (conjunção)}{}%
  \iftag{prvOr}{\ycomma $\lor$ (disjunção)}{}%
  \iftag{prvIf}{\ycomma $\lif$ (!!{conditional})}{}%
  \item O operador de conhecimento $\Knows_a$, em que $a \in G$.
  \item O operador de anúncio público $[!B]$, em que $!B$ é !!a{formula}.
\end{enumerate}
\end{defn}

O operador de anúncio público funciona como um operador caixa, e nossa
definição indutiva da linguagem é formulada de acordo com isso:

\begin{defn}
  As \emph{!!^{formula}s} da linguagem epistêmica são definidas
  indutivamente da seguinte maneira:
  \begin{enumerate}
  \tagitem{prvFalse}{$\lfalse$ é uma !!{formula} atômica.}{}

  \tagitem{prvTrue}{$\ltrue$ é uma !!{formula} atômica.}{}

  \item Toda variável proposicional $\Obj p_i$ é uma !!{formula}
  (atômica).

  \tagitem{prvNot}{Se $!A$ é !!a{formula}, então $\lnot !A$ é
  !!a{formula}.}{}

  \tagitem{prvAnd}{Se $!A$ e $!B$ são !!{formula}s, então $(!A \land
  !B)$ é !!a{formula}.}{}

  \tagitem{prvOr}{Se $!A$ e $!B$ são !!{formula}s, então $(!A \lor !B)$
  é !!a{formula}.}{}

  \tagitem{prvIf}{Se $!A$ e $!B$ são !!{formula}s, então $(!A \lif !B)$
  é !!a{formula}.}{}

  \tagitem{prvIff}{Se $!A$ e $!B$ são !!{formula}s, então $(!A \liff !B)$
  é !!a{formula}.}{}

  \item Se $!A$ é !!a{formula} e $a \in G$, então $\Knows_a !A$ é
  !!a{formula}.
  
  \item Se $!A$ e $!B$ são !!{formula}s, então $[!A] !B$ é !!a{formula}.

  \tagitem{limitClause}{Nada mais é !!a{formula}.}{}
\end{enumerate}
\end{defn}

A leitura pretendida da !!{formula} $[!A] !B$ é ``Depois que $!A$ é
anunciada de modo verídico, $!B$~vale. Às vezes também será útil falar
de conhecimento comum no contexto de anúncios públicos; por isso, a
linguagem pode incluir ainda o operador de conhecimento comum~$\CKnows_G
!A$.

\end{document}
